\section{Exponential accuracy of the aggregate event}
\label{sec:high-range}

\begin{lemma}[Exponentially accurate aggregate geometry above a logarithmic scale]
\label{restore:high-aggregate}
Fix any $A_0>0$ and $s>0$. Under the fixed positive compact structural
conditions, if ${np}\ge A_0\log n$ and $np^2\le C_1$, there is an aggregate
event $\mathcal E_s$ with
\[
 \PP(\mathcal E_s^c)\le C_s e^{-s{np}}
\]
on which the capped-adjacency geometry, weighted degree-energy bounds,
two-sided total-degree bounds, and true block-total bounds used in
\cref{restore:fit-geometry-one,restore:fit-geometry-two,restore:fit-block-event}
hold for the full graph and every graph obtained by deleting at most three
vertices. The constants may depend on $A_0,s,C_1$ and the structural
bounds. No condition is imposed on the smallest raw degree or the
largest spectral row.
\end{lemma}

\begin{proof}
The centered bounded-entry norm inequality of \citet[Corollary 3.12 and Remark 3.13]{bvh} gives
\[
 \PP\{\norm{A-{ZPZ^\top}}_{\op}>L_s\sqrt {np}\}
 \le e^{-(s+2){np}}
\]
after choosing a sufficiently large fixed $L_s$. Indeed, the row
variance sum is $O({np})$, the entry bound is one, and the prefactor $n$
in the norm-tail inequality is absorbed using $\log n\le {np}/A_0$.
The completed-mean diagonal has norm $O(p)$ and is absorbed in the
same threshold. No minimum raw degree is imposed.

The Bernoulli transform of the total number of edges gives, with
failure probability at most $Ce^{-c{n^2p}}$,
\[
 c{n^2p}\le\sum_iD_i\le C{n^2p}.
\]
For each of the fixed $K^2$ true block totals, Bernstein's inequality
for the independent upper-triangular edges, whose coefficients have
absolute value at most two, gives
\[
 \PP\left\{\left|[Z^\top(A-{ZPZ^\top})Z]_{k\ell}\right|>Ln\right\}
 \le 2\exp\left\{-\frac{cL^2n}{{np}+L}\right\}.
\]
The deterministic completed diagonal costs $O({np})$ and is included
by increasing $L$. Since ${np}\to\infty$ and $np^2\le C_1$,
the exponent is at least $c'L^2{np}/C_1$ for all sufficiently large $n$.
Increasing the fixed $L=L_s$ therefore makes the failure of all block
totals at most $C_se^{-(s+2){np}}$. The total-degree failure is also
bounded by $Ce^{-(s+2){np}}$ for large $n$.

These three full-graph events imply the required events on every
fixed-depth minor without a degree-wise probability union. To see this,
let $J=[n]\setminus T$ with $|T|\le3$, set $A_J=A_{JJ}$ and
$Z_J=Z_{J\cdot}$, and recompute all degrees and the cap
threshold on $J$. The centered matrix $A_J-{Z_JPZ_J^\top}$ is a principal
submatrix of $A-{ZPZ^\top}$, so its norm is at most $L_s\sqrt {np}$.
Removing $T$ changes the total raw degree by at most $2|T|n$;
as ${np}\to\infty$, the two-sided bound by fixed multiples of ${n^2p}$
therefore persists uniformly over $T$. Removing $T$ changes a true
block total of $A-{ZPZ^\top}$ by $O(|T|n)$ deterministically, since
the entries have bounded absolute value. Its $O_s(n)$ bound persists
as well. The surviving group masses differ from their original
values by $O(1)$ and remain comparable to $n$.

The recomputed average degree is consequently comparable to ${np}$.
For the fixed sufficiently large cap multiplier $b_0$, its threshold
$\tau_J=b_0\overline D_J$ lies between positive fixed multiples of
${np}$ and above every completed row sum $\sum_{j\prime\in J}P_{z_j z_{j\prime}}$.
Thus $G_{J,j}=\max(D_{J,j},\tau_J)\ge c{np}$ and
$\sum_jG_{J,j}\le C{n^2p}$. The raw norm bound immediately gives
\begin{align*}
 \norm{G_J^{-1/2}(A_J-{Z_JPZ_J^\top})G_J^{-1/2}}_{\op}
   &\le C_s{(np)}^{-1/2},\\
 \sum_{j\in J}\frac{(D_{J,j}-\sum_{j\prime\in J}P_{z_j z_{j\prime}})^2}{G_{J,j}}
   &=\norm{G_J^{-1/2}(A_J-{Z_JPZ_J^\top})\one}_2^2
     \le C_sn.
\end{align*}
The cap exceeds the completed population degrees, so coordinatewise
\[
 \frac{(G_{J,j}-\tau_J)^2}{G_{J,j}}
 \le\frac{(D_{J,j}-\sum_{j\prime\in J}P_{z_j z_{j\prime}})^2}{G_{J,j}}.
\]
Indeed, the left side is zero when $D_{J,j}\le\tau_J$; otherwise
$0\le D_{J,j}-\tau_J\le D_{J,j}-\sum_{j\prime\in J}P_{z_j z_{j\prime}}$. This proves the
other weighted energy bound. The order-${np}$ signal gap of the
capped adjacency matrix (a fixed gap after dividing by $\tau_J$)
and the residual
$\norm{(I-U_JU_J^\top)G_J^{1/2}Z_J}_F\le C_s\sqrt n$
now follow by the deterministic Gram-matrix and projection arguments
in \cref{prop:sparse-floor}. Those arguments use only
$G_{J,j}\ge\tau_J\asymp {np}$, total degree, the displayed rescaled
noise bound, and the energy relative to $\tau_J$, all just proved.
Intersecting the stated aggregate consequences defines
$\mathcal E_s$. The three full-graph events imply it simultaneously
for every indicated minor, proving its claimed failure bound.

For a deleted graph used as independent training data, form the
analogous event using that graph and its further minors. The same
argument gives failure probability $C_se^{-s{np}}$, and this event is
measurable from the training graph alone. The full-graph common event
is used only for deterministic implications, never as a condition
that would preserve independence of an incident row.
\end{proof}

